\clearpage
\section*{ЗАКЛЮЧЕНИЕ}
\addcontentsline{toc}{section}{ЗАКЛЮЧЕНИЕ}

Выпускная квалификационная работа была посвящена разработке и экспериментальной проверке программного решения для прогнозирования и оптимизации торговых решений на основе биржевых котировок с использованием методов обучения с подкреплением. В ходе работы была решена не только задача включения RL-алгоритмов в финансовую постановку, но и более широкая инженерная задача: создание единого программного контура, позволяющего корректно сравнивать baseline-стратегии, статистические методы, модели машинного обучения и RL-агентов.

В теоретической части были изучены особенности биржевых котировок как нестационарных финансовых временных рядов, рассмотрены классические подходы к прогнозированию, включая ARIMA/SARIMA, методы машинного обучения, нейросетевые архитектуры LSTM/GRU, а также RL-алгоритмы DQN, PPO и A2C. Отдельно был проведен обзор существующих аналогов и фреймворков — FinRL, TensorTrade и gym-anytrading. Это позволило выполнить требование о сравнении аналогов и показать, что предлагаемый подход не существует в вакууме, а занимает определенное место среди существующих исследовательских и прикладных решений.

В результате анализа была четко сформулирована проблема, которую решает работа. Обычное прогнозирование цены или направления движения недостаточно для практической оценки качества алгоритма, потому что высокая точность классификации не равна прибыльности торговой стратегии. Реальная ценность модели определяется тем, как ее сигналы влияют на динамику капитала, частоту сделок, глубину просадки и устойчивость на разных рыночных режимах. Поэтому центральным результатом работы стало построение программного стенда, в котором baseline-стратегии, ML-модели и RL-агенты сравниваются на одних и тех же данных по одинаковому набору финансовых и классификационных метрик.

В практической части реализован программный проект на Python с модульной структурой, включающей модуль загрузки данных через \texttt{yfinance}, модуль формирования OHLCV-датасета и признаков, модуль baseline-стратегий, модуль ML-моделей, торговую RL-среду на базе Gymnasium, модуль обучения агентов DQN и PPO через stable-baselines3, модуль backtesting, модуль расчета финансовых метрик и модуль визуализации результатов. Тем самым в работе не только описан теоретический подход, но и создано реальное программное решение, поддерживающее воспроизводимый запуск вычислительного эксперимента.

По итогам экспериментальной проверки разработанного ПО было показано, что:
\begin{enumerate}[leftmargin=1.25cm,label=\arabic*)]
\item стратегия Buy \& Hold сохраняет сильные позиции как простой и конкурентоспособный baseline, особенно на выраженном восходящем рынке;
\item стратегия на основе скользящих средних позволяет уменьшить просадку, но уступает пассивному удержанию по суммарной доходности на трендовом интервале;
\item среди реализованных ML-моделей наилучший баланс между качеством классификации и торговым результатом показал Gradient Boosting;
\item RL-подходы действительно способны использовать последовательную постановку задачи и оптимизировать торговые решения напрямую через функцию вознаграждения;
\item PPO оказался наиболее устойчивым из реализованных RL-агентов по совокупности риск-скорректированных метрик, тогда как DQN проявил большую чувствительность к параметрам среды и reward function.
\end{enumerate}

При этом результаты работы не дают основания утверждать, что RL всегда превосходит классические методы. Напротив, исследование показало, что качество RL-агента существенно зависит от выбора признаков, конфигурации среды, функции вознаграждения и характера рыночного режима. В ряде сценариев простые и интерпретируемые baseline-стратегии могут быть устойчивее и проще в сопровождении, а отдельные ML-модели способны демонстрировать конкурентоспособный результат при значительно меньшей вычислительной сложности.

Таким образом, основная ценность выполненной работы заключается не в декларации безусловного превосходства обучения с подкреплением, а в создании единого сравнительного программного стенда, позволяющего сопоставлять различные классы методов в корректной и воспроизводимой постановке. Разработанное решение может использоваться как основа для дальнейшего развития исследования и для расширения набора поддерживаемых торговых сценариев.

К числу перспективных направлений относятся: расширение набора признаков за счет макроэкономических и новостных данных, включение CatBoost/XGBoost, LSTM/GRU и A2C в основной экспериментальный контур, переход к многоактивному портфелю, использование walk-forward-валидации, моделирование проскальзывания и более сложных торговых ограничений, а также развитие reward function с учетом риска, просадки и стабильности стратегии.

Следовательно, цель работы достигнута: разработано, программно реализовано и экспериментально проверено решение для прогнозирования и оптимизации торговых решений на основе биржевых котировок, включающее модули загрузки данных, формирования признаков, реализации baseline- и ML-подходов, RL-среды, обучения агентов, backtesting, расчета метрик и визуализации результатов. Поставленные задачи решены, выпускная квалификационная работа выполнена в полном объеме.
